\section{Antecedentes del proyecto}
\label{sec:antecedentes}

El sistema propuesto recoge experiencias previas de digitalización de matrículas a nivel internacional, nacional y del propio Módulo San Julián, que muestran cómo pasar del registro manual a una plataforma web ordena el proceso y reduce los errores:

\subsection{Antecedentes internacionales}

En el ámbito internacional, \textcite{udb2018} analizaron en \textit{Sistema mecanizado para la inscripción de asignaturas y pagos por internet} el paso de la atención manual por ventanilla a un sistema web centralizado. En su caso, automatizar la matrícula redujo más de 80 \% el tiempo de espera y eliminó los errores al asignar vacantes, lo que mejoró la satisfacción de los estudiantes y dio más transparencia al trámite.

\textcite{sujovich2024} estudiaron la modernización de instituciones públicas con tecnologías de información y encontraron que innovar no es solo cambiar el plan de estudios, sino actualizar la forma en que se atiende a los usuarios. Observaron que muchos rechazan al inicio los sistemas digitales cuando no pueden resolver solos un problema de acceso y cuando no hay una guía clara según su nivel de manejo de la tecnología. Por eso recomendaron interfaces simples y una opción para que cada usuario recupere su acceso sin depender de otra persona.

\subsection{Antecedentes nacionales del Bachillerato Técnico Humanístico}

En Bolivia, con el Bachillerato Técnico Humanístico (BTH) y el Modelo Educativo Sociocomunitario Productivo se crearon los Módulos Tecnológicos Productivos (MTP) para compartir talleres entre varios colegios de secundaria. En Santa Cruz, por ejemplo, varios distritos han tenido el mismo problema: cuando la matrícula se hace en papel, los cupos se saturan y los horarios del Módulo chocan con los de los colegios de origen. Esa experiencia muestra que hace falta una herramienta digital que ordene la inscripción en un solo lugar.

\subsection{Antecedente institucional local: MTP San Julián}

En el ámbito específico del MÓDULO TECNOLÓGICO PRODUCTIVO SAN JULIÁN, el procedimiento de asignación de cupos atravesó dos etapas históricas determinantes que anteceden directamente a la presente propuesta:

En primer término, durante la Fase Presencial Tradicional (2022--2024), la matrícula se realizaba de manera física y manuscrita en las ventanillas del módulo. Ante la limitada capacidad de vacantes en las siete especialidades técnicas, los estudiantes y sus familias acampaban en los predios del establecimiento durante días a la intemperie, provocando desgaste físico, riesgos de seguridad y pérdidas de jornadas laborales en las labores agrícolas y comerciales de la región.

Posteriormente, en la Fase Digital Preliminar (Gestión 2025), se implementó una primera versión de software web para mitigar las aglomeraciones físicas. Aunque demostró la viabilidad del registro en línea, la plataforma evidenció tres vacíos operativos críticos: la carencia de un mecanismo de recuperación de contraseñas (lo que provocó el regreso masivo de estudiantes a pedir reseteos manuales y la pérdida de cupos en carreras de su vocación), la ausencia de un módulo para registrar y validar los requisitos documentales de la matrícula en la gestión lectiva correspondiente, y la falta de una funcionalidad para agrupar y asignar automáticamente a los estudiantes en cursos y paralelos según su colegio de origen. Esta trayectoria institucional fundamenta de manera concluyente la necesidad de la Versión 2.0 del sistema web.
